% Generated with knitr; edit this Rnw source and rebuild, not the output tex.
\section{Normal-mean interval}
The following chunk calculates the summary statistics and confidence interval
for the illustrative observations in Chapter~\ref{chap:normal-mean}.
The units are arbitrary.

\begin{knitrout}
\definecolor{shadecolor}{rgb}{0.969, 0.969, 0.969}\color{fgcolor}\begin{kframe}
\begin{alltt}
\hldef{x} \hlkwb{<-} \hlkwd{c}\hldef{(}\hlnum{1}\hldef{,} \hlnum{2}\hldef{,} \hlnum{3}\hldef{)}
\hldef{n} \hlkwb{<-} \hlkwd{length}\hldef{(x)}
\hldef{x_bar} \hlkwb{<-} \hlkwd{mean}\hldef{(x)}
\hldef{s} \hlkwb{<-} \hlkwd{sd}\hldef{(x)}
\hldef{se} \hlkwb{<-} \hldef{s} \hlopt{/} \hlkwd{sqrt}\hldef{(n)}
\hldef{ci} \hlkwb{<-} \hldef{x_bar} \hlopt{+} \hlkwd{c}\hldef{(}\hlopt{-}\hlnum{1}\hldef{,} \hlnum{1}\hldef{)} \hlopt{*} \hlkwd{qt}\hldef{(}\hlnum{0.975}\hldef{,} \hlkwc{df} \hldef{= n} \hlopt{-} \hlnum{1}\hldef{)} \hlopt{*} \hldef{se}
\hlkwd{round}\hldef{(}\hlkwd{c}\hldef{(}\hlkwc{n} \hldef{= n,} \hlkwc{mean} \hldef{= x_bar,} \hlkwc{sd} \hldef{= s,} \hlkwc{se} \hldef{= se,}
        \hlkwc{lower} \hldef{= ci[}\hlnum{1}\hldef{],} \hlkwc{upper} \hldef{= ci[}\hlnum{2}\hldef{]),} \hlnum{3}\hldef{)}
\end{alltt}
\begin{verbatim}
##      n   mean     sd     se  lower  upper 
##  3.000  2.000  1.000  0.577 -0.484  4.484
\end{verbatim}
\end{kframe}
\end{knitrout}
